\documentclass[a4paper]{article}

\usepackage{graphicx}
\usepackage{amsmath,amssymb}
\usepackage{minted}
\usepackage{multirow}
\usepackage{float}
\usepackage{algorithm}
\usepackage{algpseudocode}
\usepackage{todonotes}
\usepackage{forest}
\usepackage{xstring}
\usepackage{xcolor}
\usepackage[most]{tcolorbox}
\usepackage{xparse}
\usepackage{natbib}
\usepackage{adjustbox}

\usetikzlibrary{graphs,graphdrawing}
\usegdlibrary{layered}

% for external graphviz
\input{/home/donkarlo/Dropbox/repo/nd_language_ai_project/src/nd_language_ai/formal/computer/programming/tex/latex/code_clips/external_graphviz.tex}

% for tags
\input{/home/donkarlo/Dropbox/repo/nd_language_ai_project/src/nd_language_ai/formal/computer/programming/tex/latex/part/control_sequence/kind/semtag/semtag.tex}

% for links
\input{/home/donkarlo/Dropbox/repo/nd_language_ai_project/src/nd_language_ai/formal/computer/programming/tex/latex/code_clips/seehere.tex}

% for nested lists and sections
\input{/home/donkarlo/Dropbox/repo/nd_language_ai_project/src/nd_language_ai/formal/computer/programming/tex/latex/code_clips/deep_structure.tex}

\title{Directed Gaussian Wasserstein distance}
\date{}
\author{Mohammad Rahmani}

\begin{document}
   \maketitle

   The variables in this document correspond to the directed Gaussian Wasserstein formulation used in autonomous navigation. The Gaussian Wasserstein distance is based on the closed-form distance between multivariate normal distributions \cite{dowson-1982-the-frechet-distance-between-multivariate-normal-distributions}. Direction is added in this research from the displacement between the two Gaussian means.

   \section{Research formulation variables}

   \begin{table}[H]
      \centering
      \small
      \setlength{\tabcolsep}{4pt}
      \begin{adjustbox}{max width=\textwidth}
      \begin{tabular}{p{0.28\textwidth} p{0.20\textwidth} p{0.48\textwidth}}
         \hline
         \textbf{Concept} & \textbf{Variable name} & \textbf{Description} \\
         \hline

         Discrete time-step index
         & $k$
         & Index of the current perception--action or control step. \\

         Current state
         & $x_k$
         & Estimated current state at time step $k$. \\

         Preferred state
         & $x_k^{\mathrm{ref}}$
         & Preferred or reference state at time step $k$. \\

         Current state distribution
         & $p(x_k)$
         & Probability distribution representing the estimated current state. \\

         Preferred state distribution
         & $p(x_k^{\mathrm{ref}})$
         & Probability distribution representing the preferred or reference state. \\

         Current Gaussian mean
         & $\mu_k$
         & Mean vector of the current Gaussian state distribution. \\

         Preferred Gaussian mean
         & $\mu_k^{\mathrm{ref}}$
         & Mean vector of the preferred Gaussian state distribution. \\

         Current covariance
         & $\Sigma_k$
         & Covariance matrix of the current Gaussian state distribution. \\

         Preferred covariance
         & $\Sigma_k^{\mathrm{ref}}$
         & Covariance matrix of the preferred Gaussian state distribution. \\

         Gaussian Wasserstein distance
         & $W_2$
         & Second-order Wasserstein distance between the current and preferred Gaussian state distributions. It gives the magnitude of the distributional discrepancy. \\

         Squared Gaussian Wasserstein distance
         & $W_2^2$
         & Squared Wasserstein distance, equal to the sum of the squared mean-displacement term and the covariance-discrepancy term. \\

         Mean-displacement vector
         & $\mu_k^{\mathrm{ref}}-\mu_k$
         & Vector pointing from the current Gaussian mean toward the preferred Gaussian mean. \\

         Mean-displacement norm
         & $\left\|\mu_k^{\mathrm{ref}}-\mu_k\right\|_2$
         & Euclidean magnitude of the displacement between the current and preferred means. \\

         Unit direction toward preferred state
         & $\dfrac{\mu_k^{\mathrm{ref}}-\mu_k}{\left\|\mu_k^{\mathrm{ref}}-\mu_k\right\|_2}$
         & Unit vector that gives the direction from the current mean toward the preferred mean. \\

         Directed Wasserstein discrepancy
         & $\mathbf{w}_k$
         & Wasserstein distance multiplied by the unit mean-displacement direction. Its magnitude is $W_2$ and its direction points toward the preferred state. \\

         Position-level directed Wasserstein discrepancy
         & $\mathbf{w}_{p,k}$
         & Directed Wasserstein discrepancy expressed in physical position coordinates and passed to the outer proportional position-control loop. \\

         \hline
      \end{tabular}
      \end{adjustbox}

      \caption{Variables used in the directed Gaussian Wasserstein formulation.}
      \label{tab:directed-gaussian-wasserstein-variables}
   \end{table}

   \section{Implementation correspondence}

   \begin{table}[H]
      \centering
      \small
      \setlength{\tabcolsep}{4pt}
      \begin{adjustbox}{max width=\textwidth}
      \begin{tabular}{p{0.27\textwidth} p{0.22\textwidth} p{0.49\textwidth}}
         \hline
         \textbf{Implementation concept} & \textbf{Code variable} & \textbf{Mathematical meaning} \\
         \hline

         First Gaussian mean
         & \texttt{first\_mean}
         & Mean vector of the first Gaussian supplied to the implementation. In \texttt{calculate\_distance\_vector}, it should be the preferred mean when the desired vector direction is current $\rightarrow$ preferred. \\

         Second Gaussian mean
         & \texttt{second\_mean}
         & Mean vector of the second Gaussian. It should be the current mean when the desired vector direction is current $\rightarrow$ preferred. \\

         First Gaussian covariance
         & \texttt{first\_covariance}
         & Covariance matrix associated with \texttt{first\_mean}. \\

         Second Gaussian covariance
         & \texttt{second\_covariance}
         & Covariance matrix associated with \texttt{second\_mean}. \\

         Squared mean discrepancy
         & \texttt{mean\_distance\_squared}
         & Squared Euclidean distance between the two Gaussian means. \\

         Squared covariance discrepancy
         & \texttt{covariance\_distance\_squared}
         & Covariance contribution to the squared Gaussian Wasserstein distance. \\

         Squared total distance
         & \texttt{distance\_squared}
         & Sum of the mean and covariance contributions before taking the square root. \\

         Wasserstein distance
         & \texttt{distance}
         & Scalar $W_2$ returned by \texttt{calculate\_distance}. \\

         Direction vector
         & \texttt{direction}
         & Difference \texttt{first\_mean - second\_mean}. With preferred mean first and current mean second, this points from current toward preferred. \\

         Direction magnitude
         & \texttt{direction\_norm}
         & Euclidean norm of \texttt{direction}. \\

         Unit direction
         & \texttt{unit\_direction}
         & Normalized mean-displacement vector, \texttt{direction / direction\_norm}. \\

         Directed distance vector
         & \texttt{distance\_vector}
         & Product of scalar Wasserstein distance and unit direction; corresponds to $\mathbf{w}_k$. \\

         Numerical tolerance
         & \texttt{numerical\_tolerance}
         & Small positive threshold used to suppress floating-point errors and to detect an undefined mean-based direction when the means are effectively identical. \\

         Eigenvalues
         & \texttt{eigenvalues}
         & Eigenvalues used to compute the symmetric square root of a covariance-related matrix. \\

         Eigenvectors
         & \texttt{eigenvectors}
         & Eigenvectors used to reconstruct the symmetric matrix square root. \\

         Clipped eigenvalues
         & \texttt{clipped\_eigenvalues}
         & Eigenvalues clipped at zero to avoid small negative values caused by numerical error before taking square roots. \\

         Symmetric matrix square root
         & \texttt{square\_root\_matrix}
         & Symmetric matrix square root reconstructed from the eigenvectors and square roots of the non-negative eigenvalues. \\

         \hline
      \end{tabular}
      \end{adjustbox}

      \caption{Correspondence between the Gaussian Wasserstein implementation and the research notation.}
      \label{tab:directed-gaussian-wasserstein-implementation-variables}
   \end{table}

   \paragraph{Direction convention}
   The implementation defines
   \[
      \texttt{direction}
      =
      \texttt{first\_mean}
      -
      \texttt{second\_mean}.
   \]
   Therefore, to obtain the convention used in autonomous navigation,
   \[
      \mu_k^{\mathrm{ref}}-\mu_k,
   \]
   \texttt{first\_mean} corresponds to $\mu_k^{\mathrm{ref}}$ and \texttt{second\_mean} corresponds to $\mu_k$. The Wasserstein distance itself is symmetric with respect to the two distributions, but this added mean-based direction is order-dependent.

   \paragraph{Undefined direction}
   If the two means are equal within the numerical tolerance while the covariance matrices are different, $W_2$ can be non-zero although the mean-based direction is undefined. The implementation therefore raises an error in this case instead of assigning an arbitrary direction. If both the mean-based direction and Wasserstein distance are effectively zero, the implementation returns the zero vector.

   \bibliographystyle{plainnat}
   \bibliography{/home/donkarlo/Dropbox/repo/nd_shared_research/bibliography/bibliography.bib}
\end{document}
